\section{Related Work}

% 第一部分：从单纯的 RAG 升级到 针对复杂问题的 RAG
% 顺带提一句 Hallucination 即可，不要单列
\textbf{RAG for Complex Reasoning.}
Retrieval-Augmented Generation (RAG) serves as a robust paradigm to mitigate hallucinations by grounding Large Language Models (LLMs) in external knowledge corpora~\cite{lewis2020retrieval}. While standard dense retrieval frameworks excel at addressing simple factoid queries, they often falter in complex, multi-hop reasoning scenarios where evidence is fragmented across multiple documents~\cite{yang2018hotpotqa, thakur2021beir}. To overcome these limitations, iterative retrieval methods such as IR-CoT~\cite{trivedi2023interleaving} have been proposed to interleave reasoning traces with retrieval steps. However, these approaches primarily operate on a flat semantic space, which lack explicit modeling of the structural dependencies between entities. This limitation underscores the necessity of structured retrieval mechanisms, such as Graph-based RAG~\cite{edge2024local}, which leverage relational priors to maintain logical coherence across disparate information nodes.

\textbf{Graph-based RAG.}
Recognizing the limitations of standard RAG, recent works incorporate Knowledge Graphs (KGs) to introduce structural priors~\cite{hu2025graggraphretrievalaugmentedgeneration, mavromatis2024gnnraggraphneuralretrieval}.
Approaches like HippoRAG~\cite{jimenez2024hipporag} and its successor HippoRAG2~\cite{gutierrez2025rag} draw inspiration from hippocampal memory indexing, utilizing Personalized PageRank algorithm to navigate relational paths. By mimicking associative memory, they effectively bridge disconnected entities across the corpus. At the same time, methods such as GraphRAG~\cite{edge2024local} focus on hierarchical abstraction; they employ LLMs to pre-summarize clustered graph communities, providing a "map-reduce" style global context that handles high-level thematic queries more effectively than point-wise retrieval.
In parallel, GNN-RAG~\cite{mavromatis2024gnnraggraphneuralretrieval} introduces a neuro-symbolic framework that utilizes GNNs to reason over dense subgraphs for answer retrieval, verbalizing the extracted shortest paths into natural language to enable LLMs to perform complex multi-hop reasoning.
Despite their effectiveness in improving reasoning accuracy, these methods typically incur high computational costs and latency due to complex graph algorithms or extensive LLM usage, limiting their scalability in real-time applications.

% \subsection{Efficient Retrieval with Structural Awareness}
% Balancing structural utilization with retrieval efficiency remains a formidable challenge. Some lightweight frameworks, such as LinearRAG~\cite{zhuang2025linearrag}, attempt to accelerate graph-based retrieval by simplifying the topology. While efficient, this reductionist approach inevitably discards critical relational semantics and high-order dependencies required for precise multi-hop reasoning. Meanwhile, HypergraphRAG~\cite{luo2025hypergraphrag} introduces a hypergraph-based indexing mechanism to model non-pairwise relationships among multiple entities. By representing complex document-entity interactions as hyperedges, it facilitates more efficient information propagation across overlapping semantic clusters compared to traditional graph traversal. This high-order structural awareness allows for the capture of collective context without the prohibitive overhead of exhaustive path-searching, offering a promising direction for scalable yet semantically-rich retrieval.


\textbf{Efficient Retrieval with Structural Awareness.}
Balancing structural utilization with retrieval efficiency remains a formidable challenge. Current research primarily diverges into two extremes. On one hand, lightweight frameworks like LinearRAG~\cite{zhuang2025linearrag} achieve rapid retrieval by simplifying graph topology; however, this reductionist approach inevitably discards critical relational semantics and high-order dependencies necessary for precise multi-hop reasoning. On the other hand, structure-centric approaches prioritize reasoning depth but frequently encounter scalability bottlenecks. For instance, path-based methods such as ToG~\cite{sun2023think} rely on iterative LLM-guided beam searches, which capture complex dependencies at the cost of prohibitive inference latency. Similarly, subgraph-based methods like QA-GNN~\cite{yasunaga2021qa} and GRAG~\cite{hu2025graggraphretrievalaugmentedgeneration} try to preserve relational contexts via $k$-hop ego networks, yet incur substantial overhead during subgraph extraction and encoding. To bridge this gap, hybrid approaches~\cite{jin2024graph, jiang2025kg} use adaptive agents to dynamically select retrieval granularities (nodes, paths, or subgraphs). While these methods mitigate noise, the decision-making process of agents often introduces extensive overhead that limits real-time applicability. 
% To transcend the limitations of binary relations and computational inefficiency, HyperGraphRAG~\cite{luo2025hypergraphrag} introduces hyperedges to explicitly model complex $n$-ary relations ($n \ge 2$). This approach provides a more expressive yet efficient structural abstraction, capturing the multi-dimensional dependencies inherent in real-world knowledge without the exponential search space typical of traditional graph-based methods.